\documentclass[11pt]{article}
\usepackage[margin=1in]{geometry}
\usepackage[T1]{fontenc}
\usepackage{lmodern}
\usepackage{amsmath,amssymb,amsthm,mathtools}
\usepackage{microtype,booktabs,array,tabularx,longtable}
\usepackage{xcolor,listings,hyperref,fancyhdr}
\hypersetup{colorlinks=true,linkcolor=blue!45!black,urlcolor=blue!45!black,
  citecolor=blue!45!black,pdftitle={Unknot Recognition Workbench: Implementation Report}}
\setlength{\emergencystretch}{2em}
\setlength{\parskip}{0.35em}
\setlength{\parindent}{0pt}
\newtheorem{theorem}{Theorem}[section]
\newtheorem{proposition}[theorem]{Proposition}
\newtheorem{lemma}[theorem]{Lemma}
\newtheorem{corollary}[theorem]{Corollary}
\theoremstyle{remark}\newtheorem{remark}[theorem]{Remark}
\newcommand{\F}{\mathbb F_2}
\newcommand{\Kh}{\widetilde{\operatorname{Kh}}}
\newcommand{\code}[1]{\texttt{#1}}
\newcommand{\complexity}{n^{O(\log n)}}
\pagestyle{fancy}
\fancyhf{}\fancyhead[L]{\small Unknot recognition workbench}
\fancyhead[R]{\small Implementation report}
\fancyfoot[C]{\thepage}
\setlength{\headheight}{14pt}
\lstset{basicstyle=\ttfamily\small,columns=fullflexible,keepspaces=true,
  breaklines=true,frame=single,framerule=0.3pt,showstringspaces=false,
  aboveskip=0.8em,belowskip=0.8em}
\title{Unknot Recognition Workbench\\
\large Exact Reference Code, a Polynomial Hierarchy Subroutine,\\
and the Remaining Quasi-polynomial Obligations}
\author{Implementation and analysis prepared for Vladimir Reshetnikov}
\date{17 September 2026}
\begin{document}
\maketitle

\begin{center}
\fbox{\begin{minipage}{0.91\textwidth}
\textbf{Delivery status.} The requested general $\complexity$ implementation
has \textbf{not been completed}. The executable recognizer in this archive is
exact on completion and has exponential worst-case complexity. A genuine
polynomial-time terminal boundary-pattern routine and a conditional hierarchy
bound auditor are also implemented. Neither supplies the missing compressed
geometric hierarchy engine. No quasi-polynomial guarantee is claimed.
\end{minipage}}
\end{center}

\begin{abstract}
This report documents working, standard-library Python code rather than a
flowchart with assumed geometric oracles. The exact reference recognizer
computes reduced Khovanov homology over $\F$, with an optional determinant
obstruction. A separate routine decides essentiality of a boundary pattern on
an already known 3-ball by enumerating edge bonds of size at most three and
constructing dual-cycle witnesses. We prove its combinatorial criterion and
polynomial complexity. We also give the precise conditional arithmetic behind
$L(g+1)^L$, distinguish iteration counts from bit-operation bounds, and identify
the unimplemented constructions needed to turn the supplied talk into the
requested algorithm. The delivered test run passes 49 unittest methods. An
explicit family of unknot diagrams demonstrates exponential state growth in
the reference recognizer, so the implementation's limitation is not concealed
by small-instance timings.
\end{abstract}

\clearpage
\setcounter{tocdepth}{1}
\tableofcontents
\newpage

\section{Scope, interface, and actual deliverables}

The input is a finite classical knot diagram with $n$ crossings. The intended
answer is whether its represented knot is ambient isotopic to the unknot.
For conventional labels the diagram encoding has $O(n\log(n+2))$ bits; arbitrary
integer labels are also accepted, so parsing costs additionally depend on their
bit lengths. The requested $\complexity$ bound concerns this \emph{given}
diagram, not the minimum crossing number of its knot.

The archive has three executable components:

\begin{center}
\begin{tabularx}{\textwidth}{@{}>{\raggedright\arraybackslash}p{0.20\textwidth}>{\raggedright\arraybackslash}X >{\raggedright\arraybackslash}p{0.23\textwidth}@{}}
\toprule
Component & Implemented function & Established scope\\
\midrule
\code{recognize} & Exact determinant filter and explicit reduced Khovanov
chain complex & Exponential reference decider without ceilings\\
\code{pattern} & Essentiality on a known 3-ball; negative witness construction
and checking & Polynomial in the encoded pattern size\\
\code{bound} & Integer potential, reset checks, conditional visit count &
Arithmetic only; no topology inference\\
\bottomrule
\end{tabularx}
\end{center}

There is no recognizer backed by unspecified calls to ``find a Cheeger region''
or ``build a hierarchical multi-surface.'' There is also no hidden exponential
fallback described as having the requested bound. The reference recognizer is
explicitly a different algorithm. Its purpose is a runnable exact baseline and
small-instance test oracle for subsequent geometric development.

Run from the archive root, without installing any third-party package:
\begin{lstlisting}
python -m unknot_lab recognize examples/torus_3_5.json --check-d2
python -m unknot_lab pattern examples/pattern_prism.json
python -m unknot_lab verify-pattern examples/pattern_prism.json results/pattern_prism.json
python -m unittest discover -s tests -v
python benchmark.py
\end{lstlisting}
Python 3.10 or later is required by the package. Execution in this session used
CPython 3.13.5 on Linux; other versions were not tested. JSON results from the
recognizer explicitly contain \code{quasipolynomial\_guarantee: false}.

\section{What the supplied notes support}

The supplied document is Marc Lackenby's February 2021 talk, with 109 PDF pages,
including progressive-reveal slides \cite{notes}. The main theorem is announced
on pages 11--14. The basic hierarchy algorithm reaches its complete flowchart
on page 53: construct surfaces, cut, find a violating disc, and either obtain a
spanning disc or simplify an earlier surface and discard later ones.

The implementation-critical refinements are not merely different choices of
loop syntax. Page 74 names four speed-ups: quadratic-complexity surfaces,
compressed hierarchy encodings associated with Agol--Hass--Thurston,
multi-surfaces, and Heegaard splittings with Cheeger regions. Pages 82--91
explain how multi-surfaces are intended to shorten the effective hierarchy;
the surfaces represent independent relative homology classes and need not be
pairwise disjoint. Pages 100--108 describe the Heegaard-based size measure and
Cheeger-region alternative. The last page combines these ideas into a larger
flowchart \cite[pp.~74, 82--91, 100--109]{notes}.

This terminology matters. The talk explicitly warns that its ``genus'' can mean
an Euler-characteristic or pattern-sensitive complexity rather than ordinary
genus \cite[p.~59]{notes}. An implementation cannot substitute ordinary genus
without proving that every relevant simplification decreases the chosen
integer. Likewise, storing large multiplicities in binary does not itself
implement cutting, component extraction, surface pullback, or normalisation on
compressed objects.

\paragraph{Separate outside-source check.}
A July 2026 paper by Lackenby develops the hierarchy algorithm further.
Section 9 separates the visit count $L(g+1)^L$ from unspecified costs of some
steps. Proposition 10.2 subsequently bounds a refined hierarchy length by
$4c_q(\mathcal H)$; this is not the logarithmic-depth bound needed here. Its
Proposition 2.6 supplies a dual-graph characterization of essential patterns on
balls, which we use as an independent test criterion \cite{lackenby2026}.
These are outside findings, not additional content attributed to the attachment.

Our practical attempt therefore implements the terminal graph problem exactly,
while keeping the unimplemented geometric obligations visible in
Section~\ref{sec:missing}. This is not a claim that those obligations are
impossible to discharge.

\section{Validated diagram representations}

\subsection{Planar diagrams}
A crossing is a tuple $(a,b,c,d)$ of cyclically ordered edge labels. The
underpassing strand joins $a$ to $c$ and the overpassing strand joins $b$ to $d$.
Each label occurs exactly twice. An empty crossing list denotes precisely one
crossing-free circle. The standard trefoil fixture is
\cite{katlas}
\begin{lstlisting}
{"pd": [[1,4,2,5], [3,6,4,1], [5,2,6,3]]}
\end{lstlisting}
The implementation rejects Boolean labels, noninteger labels, missing
occurrences, and malformed crossing tuples. It first renumbers labels densely;
original labels are retained for basepoint selection and round trips.

For $n>0$ there are $4n$ half-edges, also called darts. Let $\alpha$ exchange the
two occurrences of each edge and let $\rho$ rotate the four darts at each
crossing. Face walks are the cycles of $\rho\alpha$. The underlying projection
must be connected. Its orientable rotation surface has
\[
 V=n,\qquad E=2n,\qquad \chi=V-E+F=F-n.
\]
Consequently it is a sphere precisely when $F=n+2$. This check rejects a valid
abstract rotation system on a positive-genus surface rather than silently
treating it as a classical knot diagram.

A second union--find joins opposite slots at every crossing. Its classes count
link components. Exactly one class is required. Thus the detector theorem is
never applied to an unsupported multi-component link. No unspecified
crossing-free components are appended to a nonempty diagram.

\subsection{Braid closures}
The alternative input is a strand count and a signed Artin-generator word.
The permutation induced by the word must be a single cycle. Before allocating
an array indexed by strands, the parser uses the necessary inequality
$\#\text{letters}\geq\#\text{strands}-1$; this rejects enormous, obviously
invalid strand counts supplied in a short JSON input.

Each generator creates two new lower wire ends. Positive generators place the
left upper strand over the right; negative generators reverse that choice.
After the lower ends are joined to the corresponding upper ends to close the
braid, the resulting PD passes through the same validation as an ordinary PD.
This avoids having two recognition paths with different input assumptions.

\section{The exact exponential reference recognizer}

\subsection{The chain groups and their explicit basis}
Khovanov's construction and Bar-Natan's computational exposition give the
underlying invariant \cite{khovanov,barnatan}. Our implementation discards
absolute grading shifts and the quantum grading because only total reduced
rank is needed. It does not replace homology by a graded Euler characteristic.

For each $s\in\{0,1\}^n$, smooth crossing $(a,b,c,d)$ by
\[
 s_j=0:\quad (a,b),(c,d),\qquad
 s_j=1:\quad (a,d),(b,c).
\]
Union--find computes the resulting circles as sets of original edge labels.
Choose one marked edge, hence a distinguished circle in every resolution.
Work over
\[
 A=\F[x]/(x^2).
\]
The reduced basis labels the distinguished circle by $x$; every other circle
is independently labelled by $1$ or $x$. A resolution with $k(s)$ circles
therefore has $2^{k(s)-1}$ basis elements. The group in unshifted degree $i$ is
\[
 C^i=\bigoplus_{|s|=i} \F^{\,2^{k(s)-1}},\qquad
 D_i=\dim_{\F}C^i.
\]
Basis assignments are small integers whose bits select $1$ or $x$ on the
unmarked circles. The marked circle is not a free bit.

\subsection{Differentials}
A directed cube edge changes a single zero to a one. Comparing component sets
in its source and target determines which circles are unchanged and whether
the local saddle merges or splits circles. The local maps are
\[
\begin{array}{c|cccc}
(u,v)&(1,1)&(1,x)&(x,1)&(x,x)\\\hline
m(u,v)&1&x&x&0
\end{array}
\quad\text{and}\quad
\Delta(1)=1\otimes x+x\otimes1,\qquad
\Delta(x)=x\otimes x.
\]
Unchanged circles retain their labels. These formulas preserve the condition
that the marked circle has label $x$. If a marked circle merges with another
$x$-labelled circle, the image is zero, not an invalid reduced basis element.

\begin{proposition}
The implemented maps satisfy $d^{i+1}d^i=0$.
\end{proposition}
\begin{proof}
A two-step path through the cube changes two distinct crossings. The two orders
of changing them give the two sides of a square. For disjoint local saddles the
maps act on different factors and commute. In the interacting cases the needed
identities are associativity of $m$, coassociativity of $\Delta$, and
\[
 (m\otimes1)(1\otimes\Delta)=\Delta m
 =(1\otimes m)(\Delta\otimes1),
\]
together with commutativity and cocommutativity. Each identity follows directly
by evaluating its two sides on the $1,x$ basis. Thus both paths through every
square induce the same map. Their sum is zero over $\F$. Every term in $d^2$
appears in exactly such a pair, and restriction to the reduced subcomplex
preserves the identity.
\end{proof}

No cube signs are needed over $\F$. For comparison with the integral invariant,
the ordinary signed integral cube reduces to exactly these maps modulo two.
The optional \code{--check-d2} forms differential columns and composes successive
matrices. This is an additional execution-time consistency test, independent
of the rank computation.

\subsection{Bit-packed exact linear algebra}
Each differential column is a Python integer. A set bit specifies a target
basis vector. Terms cancel by XOR. Gaussian elimination selects the highest
set bit as pivot; a dictionary stores one reduced column for each pivot.
No floating-point arithmetic or numerical rank threshold is used.

Writing $r_i=\operatorname{rank}(d^i)$ and $r_{-1}=r_n=0$, the program computes
\[
 h_i=D_i-r_{i-1}-r_i,\qquad H=\sum_{i=0}^{n}h_i.
\]
Indeed, $\dim\ker d^i=D_i-r_i$, and the $r_{i-1}$-dimensional previous image is
contained in that kernel. Negative dimensions or total rank zero trigger an
internal-error result, never a knot verdict.

\subsection{Why rank one is a complete decision criterion}
The external topological input is Kronheimer--Mrowka's rank-one theorem for
reduced Khovanov homology \cite{kronheimermrowka}. The use of $\F$, rather than
rational coefficients, can be justified without an additional detection
conjecture.

\begin{proposition}
For a classical knot $K$, $\dim_{\F}\Kh(K;\F)=1$ if and only if $K$ is the
unknot.
\end{proposition}
\begin{proof}
Take the signed reduced integral complex. Its groups are free abelian, so
universal coefficients imply, degree by degree, that the homology dimension
modulo two is at least the corresponding rational dimension. The total
rational reduced homology is nonzero: otherwise its graded Euler
characteristic, the reduced Jones invariant, would vanish, contrary to its
normalization at $1$. Thus total dimension one modulo two implies total
rational dimension one. The rank-one theorem gives the unknot. Conversely,
the crossing-free unknot has one reduced basis element and zero differential;
invariance gives dimension one for every diagram of that knot.
\end{proof}

The deep detector theorem is cited, not reproved or claimed to have been
formalized. The code-level argument identifies the implemented complex with
the usual construction, but the program has not been checked in a proof
assistant.

\subsection{A sound optional determinant filter}
The knot determinant is the absolute Alexander evaluation at $-1$
\cite{determinant}. The program constructs the Fox-colouring relation matrix:
each row has $2$ in the over-arc column and $-1$ in each under-arc column, with
coefficients combined when columns coincide. A cofactor is evaluated using
fraction-free Bareiss elimination and exact integer division.

A determinant different from one proves nontriviality. Determinant one does
\emph{not} prove triviality: it causes the program to continue to homology.
The delivered $T(3,5)$ example exercises exactly this branch. Its observed
values are determinant one and reduced $\F$ rank seven.

\section{Running time: an explicit exponential lower bound}

Let $D=\sum_iD_i$. A smoothing of a connected $n$-crossing projection has at
most $n+1$ circles. To see this, join its circles at the original crossing
locations: the resulting graph is connected and has at most $n$ edges, so it
has at most $n+1$ vertices. Therefore
\[
 D=\sum_s2^{k(s)-1}\leq 2^n2^n=4^n.
\]
All saddle construction is polynomial in $n$ per state or cube edge. Elimination
on bit vectors of length at most $D$ has a coarse $O(D^3)$ bit-operation bound,
up to polynomial factors in $n$. Storing pivots takes at most $O(D^2)$ bits,
again up to polynomial bookkeeping factors. Hence the unbounded reference
implementation has a coarse $2^{O(n)}$ time and space bound, plus polynomial
input-reading costs. Treating a Python integer as a constant-cost object would
be incorrect; its vector bit length is included here.

\begin{proposition}\label{prop:lower}
The unbounded reference implementation does not have a $\complexity$
worst-case running-time bound, even with the determinant filter enabled.
\end{proposition}
\begin{proof}
Consider the closure of
\[
 \beta_n=\sigma_1\sigma_2\cdots\sigma_n\in B_{n+1}.
\]
Each last strand admits a Markov destabilization, so the closure is an unknot
with $n$ displayed crossings. Its determinant is one, and the filter cannot
finish the computation. The cube constructor explicitly visits all $2^n$
states. This is already an $\Omega(2^n)$ operation lower bound, which exceeds
$n^{C\log n}$ for every fixed $C$ at sufficiently large $n$.

More precisely, in this family's input convention a state of weight $i$ has
$i+1$ circles. This follows by removing the last single-crossing strand: one
smoothing leaves the smaller diagram, the other also splits off a circle.
Thus
\[
 D_i=\binom ni2^i,\qquad D=(1+2)^n=3^n.
\]
The growth is a feature of the specified implementation, not a conjecture about
all possible unknot algorithms.
\end{proof}

The benchmark checks these exact state and generator counts for $1\leq n\leq8$.
The proof, not those finitely many timings, establishes the asymptotic
limitation. With default ceilings the program instead eventually returns
\code{unknown}; a time- or size-limited partial procedure is not a complete
quasi-polynomial decider.

\section{A polynomial terminal boundary-pattern algorithm}
\label{sec:pattern}

\subsection{Domain and encoding}
The notes define a boundary pattern $P$ as embedded trivalent graphs and simple
closed curves. On a ball it is essential when every disc boundary meeting $P$
transversely at most three times bounds a boundary disc whose pattern is empty,
a single arc, or a tripod \cite[pp.~24--25]{notes}.

\textbf{The ambient manifold is assumed already to be a 3-ball.} A vanishing
first Betti number or a spherical boundary alone is not the required input
certificate for an arbitrary manifold. The present code does not infer this
assumption from such data.

Each graph vertex records a cyclic order of three half-edge IDs. IDs $2e$ and
$2e+1$ form edge $e$. Every ID must occur exactly once. Loops and parallel edges
are permitted. The rotation surface of each component is checked to be a
sphere. Components with no vertices are recorded by an integer circle count.
For this particular decision problem, nesting of distinct components is
irrelevant; the code does not claim to reconstruct that nesting.

\subsection{The edge-cut criterion}
For a connected cubic plane multigraph $G$, write $\delta(W)$ for the edges
with one endpoint in a nonempty proper vertex subset $W$. A \emph{bond} is a
cut for which both shores induce connected subgraphs. A three-edge bond is
\emph{trivial} when one shore is a single vertex.

\begin{theorem}\label{thm:pattern}
A boundary pattern on a 3-ball is essential precisely in the following cases:
it is empty; it is one vertex-free circle; or it is a connected cubic plane
multigraph with no one-edge bond, no two-edge bond, and no nontrivial
three-edge bond.
\end{theorem}
\begin{proof}
If the pattern has at least two connected components, a simple curve in their
complement can separate some components from the others. Neither complementary
disc has empty pattern, and the curve has no intersections with the pattern.
Pushing the curve's spanning disc into the ball gives a violating disc. Empty
pattern and one circle satisfy the definition directly.

Suppose now that the pattern is a connected cubic graph. A bond in a plane
graph corresponds to a simple cycle in its geometric dual: boundary arcs in
the dual faces trace a separating curve that crosses precisely the bond edges.
For a bond of size one no allowed disc pattern has that boundary-intersection
count. For size two both shores contain vertices, so neither side is a single
arc or empty. For a nontrivial bond of size three, neither side consists of
one vertex and three exiting edges, so neither side is a tripod. These bonds
therefore produce violations.

Conversely, suppose all three forbidden bond types are absent. Let $C$ be a
simple curve transverse to the pattern with at most three intersections. If
one complementary disc contains no graph vertex, it contains at most one arc,
or is empty. There is no detached circle in a connected cubic graph. Such a
curve is not violating.

It remains that both discs contain vertices. Let $W$ be the vertices on one
side. Every edge of $\delta(W)$ crosses $C$, so $|\delta(W)|\leq3$. Every
nonempty edge cut contains a bond. Absence of one- and two-edge bonds therefore
implies that every nonempty cut has size at least three. Consequently
$|\delta(W)|=3$, and each cut edge crosses $C$ exactly once.

Both shores must be connected: if one had at least two components, each
component would have an edge cut of size at least three, forcing at least six
edges across that shore's boundary. Thus $\delta(W)$ is a three-edge bond.
It must be trivial. Its singleton shore is a cubic vertex with three exiting
edges and no internal loop, giving exactly an allowed tripod. Hence no
violating curve exists.
\end{proof}

The dual bond--cycle correspondence used here can also be seen by taking a
small regular neighbourhood of a connected shore. Because the complementary
shore is connected, the boundary component separating the two shores traces
one dual cycle. A bridge yields a dual loop; a two-edge bond gives a dual
two-cycle. These cases must not be removed by treating the dual as an ordinary
simple graph before constructing witnesses.

\subsection{Implementation and cost}
The production algorithm examines all sets of one, two, then three edges.
For each set it deletes those edges and computes connected components. Exactly
two components with that actual cut define a bond. Size-three cuts with a
singleton smaller shore are discarded as allowable tripods.

There are $O(E^3)$ candidate sets, and a connectivity scan costs $O(V+E)$.
The decision procedure therefore costs
\[
 O\bigl(E^3(V+E)\bigr)
\]
combinatorial operations for a nonempty cubic graph. The encoding and integer
indices introduce at most polynomial bit overhead. Empty and circle-only
patterns are handled directly without expanding the circle count.

A negative graph certificate records a shore, its cut-edge IDs, and the ordered
edges and faces of the dual cycle. The verifier recomputes the cut and checks
connectivity and the cycle, without repeating the edge-subset search. For a
disconnected pattern the witness is the verified component count: it certifies
existence of a zero-intersection separating curve, not explicit geometric
coordinates for one. An essential answer follows from exhaustive absence of
small forbidden cuts, not from a claimed short positive certificate.

The tests independently form the dual and use the simple four-connected
criterion, including the complete graphs on three and four vertices,
from \cite[Proposition~2.6]{lackenby2026}. This dual-vertex search differs from
the production primal-edge search. All tested spherical rotations give the
same answer.

\section{The hierarchy potential: a conditional executable audit}
\label{sec:bound}

The notes' informal estimate treats surface complexities as bounded digits
\cite[pp.~60--65]{notes}. Here we state a precise arithmetic variant. Assume
that complete phase states carry $L$ nonnegative digits $d_1,\ldots,d_L$, each
at most $g$, with shorter lists padded by zero. Assume that a reset preserves
all digits before some pivot $j$ and strictly lowers $d_j$, while allowing
arbitrary later digits within the same bound.

Put $B=g+1$ and define
\[
 \Phi(d_1,\ldots,d_L)=\sum_{i=1}^L d_iB^{L-i}.
\]
This explicit radix handles the endpoint digit $g$; the talk uses informal
base-$g$ notation. We do not silently interpret its unspecified ``genus'' as
the concrete integer used by this module.

\begin{lemma}
Every reset lowers $\Phi$ by at least one. The number of complete phases is at
most $(g+1)^L$. If each phase makes at most $L$ loop visits, the visit count is
at most $L(g+1)^L$.
\end{lemma}
\begin{proof}
For $g\geq1$, a one-unit reduction of the pivot lowers the leading contribution
by $B^{L-j}$. The largest possible increase from all later digits is
\[
 g\sum_{k=0}^{L-j-1}B^k=B^{L-j}-1.
\]
The net reduction is therefore at least one. Also
$0\leq\Phi\leq B^L-1$, so at most $B^L$ phase states can be visited in strict
descending order. Multiply by the visits per phase. If $g=0$, all digits are
zero, no reset is possible, and the single-phase bound remains valid.
\end{proof}

A hierarchy-building step that appends a surface can increase $\Phi$. The
comparison must be between complete phases surrounding a simplification and
rebuild, not every individual instruction. The code's \code{verify\_reset}
checks the shared prefix, strict pivot decrease, digit bounds, and actual
potential decrease. These are integer assertions, not computed topological
invariants.

\begin{theorem}[One sufficient quasi-polynomial contract]
Suppose an implemented geometric engine has uniform bounds
\[
 g(n)\leq C(n+2)^a,\quad L(n)\leq b\log_2(n+2)+c,
\]
its full-phase resets satisfy the preceding hypotheses, it has at most
$R(n)$ outer restarts, and every visit costs at most $P(n)$ bit operations.
If $R$ and $P$ are polynomially bounded, its total running time is
$(n+2)^{O(\log(n+2))}$.
\end{theorem}
\begin{proof}
The total cost is bounded, up to constant control-flow overhead, by
\[
 P(n)R(n)L(n)(g(n)+1)^{L(n)}.
\]
Its logarithm is $O(\log n)+O(\log n)O(\log n)=O((\log n)^2)$.
Exponentiating gives the claim. All bounds must hold uniformly for actual
stored states, including states after restarts.
\end{proof}

These are sufficient hypotheses, not claims established for the delivered
reference code. The theorem would also admit suitable quasi-polynomial
per-visit or restart bounds, but merely asserting them does not implement
anything. Likewise, a polynomial number of bits in a surface's normal vector
does not establish a polynomial bound on operations that expand its discs.

The function \code{cheeger\_inequality\_only} evaluates
\[
 3g(\partial M_i)\leq
 \min\{g(H_i),\ g(H)-g(H_i)\}.
\]
The notes additionally require the restricted Morse function to be a Heegaard
Morse function \cite[p.~108]{notes}. The numeric inequality alone neither
recognizes such a region nor constructs the resulting modification.

\section{Concrete remaining implementation obligations}
\label{sec:missing}

The following are missing geometric constructions, not optional performance
polish on existing code. Naming their outputs is not a substitute for
constructing them.

\paragraph{Initial layered handle structure.}
Convert the diagram to the exterior with an explicit embedding model and
layered handle/Morse data. Retain boundary identifications and peripheral
information needed to recognize a compression disc for the original boundary.
Establish the initial size and bit-length bounds. The current PD parser does
not create this structure.

\paragraph{Compressed cuts and surface transport.}
Specify local cell types, attaching maps, multiplicity encodings, and interval
pairings. Implement cutting and re-gluing, Euler-characteristic and component
calculations, orientation, and maps back to earlier stages without expanding
exponentially many parallel discs. Any appeal to compressed orbit counting
must include the actual encoded operations and their cost bounds.

\paragraph{Hierarchical multi-surfaces.}
Construct the relative homology classes and their embedded representatives,
including intersection data and the sequential cuts that convert them into
hierarchy stages. Prove independence and the required complexity control.
Compression must update the multi-surface representation and identify precisely
which later data remain valid. These tasks correspond to named boxes on the
last slide, but there is no implementation of them in this archive.

\paragraph{Geometric Cheeger regions and weak reduction.}
Determine the relevant level surfaces and their intersections with each
submanifold. Verify the Morse/Heegaard hypotheses, not just an integer
inequality. Construct the new generalized splitting, implement weak reduction
and the use of Haken's lemma, and bound all resulting restarts and changes of
representation.

\paragraph{Terminal topology and witness pullback.}
Establish the hypotheses permitting the terminal regions to be treated as
balls. Apply the implemented boundary-pattern subroutine there. Transport its
disc witness through the hierarchy to either a genuine original-boundary
compression or a justified earlier simplification. A dual-cycle witness on an
abstract ball pattern by itself is not a knot spanning-disc certificate.

\paragraph{End-to-end complexity proof.}
Prove the appropriate global digit and depth bounds for the implemented
objects, including the pattern-sensitive meaning of the digits. Account for
normalization, large-integer arithmetic, compressed surface updates, and
outer restarts. Only after these obligations are met would the conditional
bound in Section~\ref{sec:bound} become a runtime theorem for an actual
quasi-polynomial recognizer.

\section{Testing, recorded results, and operational limits}

The delivered \code{results/unittest.txt} records 49 passing test methods.
Methods contain additional parameterized cases. The linear-algebra tests
compare all $512$ binary $3\times3$ matrices with a separate dense elimination
routine, use $100$ seeded random binary matrices, and compare $140$ integer
determinants of orders zero through six with rational Gaussian elimination.

Knot tests cover mirrors, every basepoint of the figure-eight example,
relabeling, crossing reordering, half-turns of crossing tuples, inverse braid
pairs, the braid relation, positive and negative stabilizations, and 24 seeded
conjugates of unknot braid words. Homology computations in these tests check
$d^2=0$. Pattern tests include all spherical choices of local rotations of
several fixed multigraphs, independent dual-criterion comparisons, and
rejection of modified certificates. CLI tests check invalid inputs, resource
limits, exact verdicts, and the standalone witness verifier.

\begin{center}
\begin{tabular}{@{}lrrrr@{}}
\toprule
Example & Crossings & Determinant & Reduced $\F$ rank & Verdict\\
\midrule
Crossing-free circle & 0 & 1 & 1 & Unknot\\
One curl & 1 & 1 & 1 & Unknot\\
Trefoil & 3 & 3 & 3 & Nontrivial\\
Figure-eight & 4 & 5 & 5 & Nontrivial\\
$T(3,5)$ & 10 & 1 & 7 & Nontrivial\\
Conjugated unknot fixture & 8 & 1 & 1 & Unknot\\
\bottomrule
\end{tabular}
\end{center}
These are recorded computations, not new knot-theoretic claims. Several ranks
in the unit suite are fixed regression expectations. No independent complete
Khovanov implementation was executed, and these tests do not formally verify
the recognizer. The different dense rank and dual-pattern algorithms provide
independent checks of their respective components, not of the entire program.

\begin{center}
\begin{tabular}{@{}rrrr@{}}
\toprule
$n$ in Proposition~\ref{prop:lower} & Resolutions & Chain generators & Logical matrix bits\\
\midrule
1 & 2 & 3 & 2\\
2 & 4 & 9 & 20\\
3 & 8 & 27 & 174\\
4 & 16 & 81 & 1,480\\
5 & 32 & 243 & 12,570\\
6 & 64 & 729 & 107,100\\
7 & 128 & 2,187 & 916,230\\
8 & 256 & 6,561 & 7,869,200\\
\bottomrule
\end{tabular}
\end{center}
Timing and environment data are in \code{results/validation.json};
\code{benchmark.py} regenerates them. Timings are not used as complexity evidence.

\paragraph{Resource and error semantics.}
The default limits are 65,536 states, 200,000 generators, and 128,000,000
logical matrix bits, where the last quantity is
$\sum_iD_iD_{i+1}$. It is a conservative combinatorial estimate, not a Python
memory-usage bound. The program does not enforce a wall-clock or RSS limit.
The \code{--unlimited} option removes all three ceilings and can be expensive.
A stopped computation returns \code{unknown} with a null Boolean, never
\code{nontrivial}. Exact yes and no answers both exit with code zero; resource
stops exit two, invalid data three, and detected internal arithmetic errors
four. Argparse also uses exit two for command syntax errors, with a message on
standard error rather than a JSON mathematical result.

\section{Conclusion}
\enlargethispage{2\baselineskip}

There is working exact recognition code, a proved polynomial terminal
boundary-pattern algorithm, a conditional complexity auditor, examples, and a
reproducible test suite in this archive. The unbounded recognizer is a complete
exponential decision procedure, not a realization of the talk's announced
quasi-polynomial algorithm. Proposition~\ref{prop:lower} explicitly rules out
that claim for this implementation. The missing geometric constructions and
uniform bit-cost bounds are listed above so that subsequent work can address
the actual requirements rather than inherit an unjustified runtime label.

\begin{thebibliography}{9}
\raggedright
\bibitem{notes}
Marc Lackenby. \emph{Unknot recognition in quasi-polynomial time}.
February 2021 talk slides. User-supplied file
\code{quasipolynomial-talk.pdf}, 109 PDF pages. Page numbers in this report refer
to those PDF pages, including reveal slides. Not redistributed in the archive.

\bibitem{lackenby2026}
Marc Lackenby. \emph{Incompressible surfaces, hierarchies and unknot recognition}.
arXiv:2607.23350v1, 25 July 2026.
\url{https://arxiv.org/html/2607.23350v1}.

\bibitem{khovanov}
Mikhail Khovanov. \emph{A categorification of the Jones polynomial}.
arXiv:math/9908171.
\url{https://arxiv.org/abs/math/9908171}.

\bibitem{barnatan}
Dror Bar-Natan. \emph{On Khovanov's categorification of the Jones polynomial}.
Publication and computational exposition page.
\url{https://www.math.toronto.edu/~drorbn/papers/Categorification/}.

\bibitem{kronheimermrowka}
Peter B. Kronheimer and Tomasz S. Mrowka.
\emph{Khovanov homology is an unknot-detector}. arXiv:1005.4346.
\url{https://arxiv.org/abs/1005.4346}.

\bibitem{katlas}
Knot Atlas. \emph{The trefoil $3_1$}. Standard planar-diagram fixture.
\url{https://katlas.org/wiki/3_1}.

\bibitem{determinant}
Knot Atlas. \emph{The Alexander--Conway Polynomial}. Determinant convention.
\url{https://katlas.org/wiki/The_Alexander-Conway_Polynomial}.
\end{thebibliography}
\end{document}
